\chapter*{How to Use This Guide}\label{ch:howto}
\addcontentsline{toc}{chapter}{How to Use This Guide}
\markboth{How to Use This Guide}{How to Use This Guide}

\section*{The project in sixty seconds}

Your thesis showed that plain language can control a real-time speech enhancer.
A sentence such as \enquote{the background noise is unbearable} becomes an
intent class and a severity score, and those become DSP parameter updates that
take effect without interrupting the audio. It ran on a laptop, in 64-bit
floating point, with a 110-million-parameter language model.

This project asks the hardware version of the same question: \emph{what is the
smallest, lowest-power hardware that can still do intent-controlled speech
enhancement, and what does each design decision cost?} You will answer it by:
\begin{enumerate}
  \item turning the thesis code into a bit-exact fixed-point model;
  \item designing the audio datapath in SystemVerilog;
  \item verifying it against that model;
  \item running it in real time on an FPGA, with your NLP pipeline controlling it;
  \item going deeper on either ML acceleration or the ASIC flow.
\end{enumerate}
Along the way you will practise the core skills of VLSI-based DSP and ML
hardware: fixed-point design, microarchitecture, RTL, verification, FPGA
prototyping, and measuring area, power and latency.

\section*{How the guide is organised}

\begin{description}[leftmargin=0pt,style=nextline]
  \item[Part I --- Foundations (read first)]
  \Cref{ch:overview} sets out the research questions, the plan and the
  deliverables. \Cref{ch:primer} is a primer on the hardware concepts you will
  need. Skim it now and come back to each section when a phase uses it.
  \Cref{ch:workload} measures your thesis system as a hardware workload. It is
  the chapter that makes everything else concrete, so do not skip it.
  \item[Part II --- The project, step by step]
  One chapter per phase, \cref{ch:setup} to \cref{ch:asic}. Every phase chapter
  has the same shape:
  \begin{enumerate}
    \item a goal;
    \item the inputs you need and the outputs you will produce;
    \item a sequence of numbered \emph{steps}, each saying what to do, how to do
          it, how you know it worked, and roughly how long it takes;
    \item the pitfalls;
    \item exit criteria and a checklist.
  \end{enumerate}
  \item[Part III --- Doing it well]
  Evaluation methodology (\cref{ch:evaluation}), writing up and presenting
  (\cref{ch:writeup}), and the timeline, milestones and risk register
  (\cref{ch:plan}).
  \item[Appendices]
  The measured fact sheet (\cref{app:factsheet}), the hardware register map
  (\cref{app:registers}), a command cheat sheet (\cref{app:cheatsheet}), a
  glossary (\cref{app:glossary}), copy-paste templates for your lab notebook and
  reports (\cref{app:templates}), and an annotated reading list
  (\cref{app:reading}).
\end{description}

\section*{Suggested reading paths}

\begin{itemize}
  \item \textbf{\enquote{I want to start tomorrow.}} Read
        \cref{ch:overview,ch:workload,ch:plan}, then begin \cref{ch:setup}.
  \item \textbf{\enquote{Hardware is new to me.}} Add \cref{ch:primer}, one
        section per evening, alongside Phase~0. Keep \cref{app:glossary} open.
  \item \textbf{\enquote{I have less time than the full plan assumes.}}
        \Cref{ch:plan} has cut-down versions of the project for different
        credit loads. Decide your scope before you start Phase~1.
\end{itemize}

\section*{The boxes}

The guide uses a small set of boxes consistently. This is what each one means:

\begin{stepbox}[title={\faShoePrints\enspace Step~N.M\quad An action for you to take}]{}
Numbered per chapter. It says what to do, how to do it, how to check it worked,
and how long it usually takes. Steps are meant to be done in order.
\end{stepbox}
\begin{donebox}
Exit criteria. When everything here is true, the phase is finished and its
milestone is reached.
\end{donebox}
\begin{whybox}
Connects a step or concept to your thesis and to the research questions, so you
always know why you are doing something.
\end{whybox}
\begin{conceptbox}{Background theory}
A compact explanation of a concept, usually with the one equation that matters.
\end{conceptbox}
\begin{rqbox}{Research question or hypothesis}
A question the project answers, or a claim it tests.
\end{rqbox}
\begin{tipbox}
Practical advice that saves time.
\end{tipbox}
\begin{pitfallbox}
A mistake people really make at this point, and how to avoid it.
\end{pitfallbox}
\begin{readbox}
Where to read more. Full references are in the bibliography and in
\cref{app:reading}.
\end{readbox}

Checklists look like this; tick them in your printed copy or in your notebook:
\begin{checklist}
  \item Read \cref{ch:overview}.
  \item Reproduce the measurements in \cref{ch:workload}.
\end{checklist}

\section*{Conventions}

\begin{itemize}
  \item \textbf{Measured or estimated.} Any number about the thesis system is
        either \emph{measured} from the thesis code or marked as an
        \emph{estimate}. Never quote an estimate as a result in your own report.
        Measure it.
  \item \textbf{Paths.} \file{src/intent_se/...} refers to the thesis repository.
        Paths such as \file{golden/fixedpoint.py} or \file{rtl/bin_pipeline.sv}
        refer to the new project repository, \file{intent-se-hw}, which you
        create in \cref{ch:setup}.
  \item \textbf{Weeks and effort.} Week markers such as \weeks{3--7} assume about
        12--15 hours of work per week. Effort estimates assume you are new to
        the tool in question. Expect the first step with any tool to take longer
        than estimated and later steps to take less.
  \item \textbf{Tools change.} Commands and versions were checked when the guide
        was written, but tools move fast. When a command fails, read the tool's
        current documentation before assuming you made a mistake. Record the tool
        versions you actually used, as \cref{ch:evaluation} explains.
\end{itemize}

\section*{This guide is a living document}

The guide is a LaTeX project inside your thesis repository
(\file{docs/research_guide/}). Build it with \code{make} in that directory (see
the \file{README.md} there). Edit it as the project evolves: when you change the
plan, change the guide, and note why in your lab notebook. A plan you keep up to
date is worth far more than a perfect one you stop reading.
